\centerline{\bf \Large Экзаменационная работа по алгебре}
\centerline{\bf \Large 9"в"\ физико-математический класс}
\centerline{\bf \Large Зимняя сессия}
\centerline{\bf \Large Вариант Q}

\begin{flushright}
        \scriptsize{Если математика это язык, на котором говорят все точные науки,\\то алгебра --- алфавит этого языка.\\РВ}
\end{flushright}

\problem{1}
Решить уравнение 

$\dfrac{\left(x^2+18 x+45\right)^2}{5}+\dfrac{5\left(2 x^2+\dfrac{14}{5} x-\dfrac{48}{5}\right)^2}{4}=\dfrac{\left(x^2+20 x+51\right)^2}{5}$

\problem{2}
Решите неравенство 

$\dfrac{\sqrt{\left(x^2-81\right)\left(x^2+10 x+25\right)}}{\left(169-x^2\right)\left(x^2+12 x+11\right)} \leqslant 0$.


\problem{3}
Постройте график функции $y = \dfrac{x^2 - 6x + 5}{x - |x-2|}$ и укажите её область. Как значения функция принимает

а) два раза \qquad\qquad\qquad б) три раза

\problem{4}
Изобразите на координатной плоскости множество точек ($x ; y$), удовлетворяющих системе:

$$
\left\{\begin{array}{l}
r^2 \leq x^2+y^2 \leq 9 r^2, \\
y-3 x \leq 0 \\
3 y+x \geq 0
\end{array}\right.
$$


При каких $r>0$ площадь полученной фигуры равна $18 \pi ?$

\problem{5}
Найдите все значения параметра $a$, при каждом из которых неравенство

$$
\left(x^2-(a+8) x-6 a^2+24 a\right) \sqrt{3-x} \leqslant 0
$$

имеет единственное решение.


\hspace{3mm}

\centerline{\bf \large МБОУ "Элистинский Лицей"}


\newpage

\hspace{3mm}

\centerline{\bf \Large Экзаменационная работа по алгебре}
\centerline{\bf \Large 9"в"\ физико-математический класс}
\centerline{\bf \Large Зимняя сессия}
\centerline{\bf \Large Вариант Q}

\begin{flushright}
        \scriptsize{Если математика это язык, на котором говорят все точные науки,\\то алгебра --- алфавит этого языка.\\РВ}
\end{flushright}

\problem{6}
Из пунктов $A$ и $B$, расположенных ша расстоянии $50$ км, на встречу друг другу одновременно вышли два пешехода. Через 5 ч они встретились. После встречи скорость первого пешехода, идущего из $A$ в $B$, уменьшилась на $1$ км/ч, а скорость второго возросла на $1$ км/ч. Найдите первоначальную скорость первого пешехода, если он прибыл в пункт $B$ на 2 ч раньше, чем второй а пункт $A$.

\problem{7}
Анджур с друзьями решили устроить пикник. Для этого им от пункта А нужно добраться вниз по реке до пункта В, причем в их распоряжении есть два катера. Считая себя самым ответственным, Анджур вызвался самостоятельно доехать до пункта В на более быстроходном катере и начать готовить место для пикника. Оба катера вышли одновременно из пункта A. Однако, промчавшись восемь километров, Анджур заметил на берегу машущего ему рукой Расула, который просил по старой дружбе довезти его до пункта С. И хоть пункт С Анджур уже проехал, он согласился. По пути в пункт С Анджур с Расулом встретили идущий навстречу второй катер с друзьями Анджура, откуда те крикнули, что им до пункта В осталась треть пути и чтобы Анджур нигде не задерживался. Доставив Расула в пункт С, Анджур немедленно помчался догонять друзей. Найдите расстояние между пунктами В и С, если известно, что оба катера пришли в пункт В одновременно, скорости катеров постоянны, а Анджур, действительно, нигде не задерживался.

\problem{8$^{*}$}
\textbf{Удивительная задача от ЭлМШ.} Аркадий Владимирович написал на доске уравнение вида  $x^2 + px + q = 0$  с целыми ненулевыми коэффициентами $p$ и $q$. Временами к доске подходили ученики 9в класса, стирали уравнение, после чего составляли и записывали уравнение такого же вида, корнями которого являются коэффициенты стёртого уравнения. В какой-то момент составленное уравнение совпало с тем, что было написано на доске изначально. Какое уравнение изначально было написано на доске?

\vspace{-0.5mm}

\centerline{\bf \large МБОУ "Элистинский Лицей"}